(1) 两光场相加再作时间平均，得 $I(\phi)=I_1+I_2+2\mathcal V\sqrt{I_1I_2}\cos\phi$，其中 $0\le\mathcal V\le1$。探测到的条纹可见度为 $(I_{\max}-I_{\min})/(I_{\max}+I_{\min})=2\mathcal V\sqrt{I_1I_2}/(I_1+I_2)$，即使光源完全相干也可能因强度不等降低对比度。
(2) 镜面位移 $\Delta x$ 使光在该臂的往返路程变化 $2\Delta x$，因此 $\Delta\phi=(2\pi/\lambda)(2\Delta x)=4\pi\Delta x/\lambda$。条纹经过的有符号次数 $N=\Delta\phi/(2\pi)=2\Delta x/\lambda$；计数总条纹时取其绝对值。
(3) 平板替代相同厚度空气的单程额外光程为 $(n-1)d$；往返路径穿过两次，故 $\Delta\mathrm{OPD}=2(n-1)d$，相位增量为 $\Delta\phi_{\rm plate}=4\pi(n-1)d/\lambda$。若厚度或折射率随波长变化，应使用实际光程而不能只用几何距离。
(4) 镜面位移灵敏度是 $d\phi/dx=4\pi/\lambda$，所以在其他参数可控时一条完整条纹对应镜面移动 $\lambda/2$。不等臂强度通过因子 $2\sqrt{I_1I_2}/(I_1+I_2)$ 降低条纹可见度，而部分相干进一步通过 $\mathcal V$ 降低它；二者改变测量噪声，却不改变位移到几何相位的基本系数。
